\section{Statistical Methodology}
\label{sec:stats}

A leaderboard that shows bare numbers violates the field's first rule. With few
scenarios, model orderings are dominated by sampling noise rather than real
capability differences, and the honest board has to say so.

\subsection{Scenario-level paired bootstrap}

Uncertainty is estimated with the nonparametric bootstrap: $B = 2000$ replicates,
fixed seed 42. The unit of resampling is the \textbf{scenario}, not the individual
row. All runs of a scenario share the same task and persona, so their scores are
correlated; resampling rows would treat correlated repeats as independent draws and
understate uncertainty. Each replicate draws scenario IDs with replacement (keeping
every run of a resampled scenario) and recomputes the per-model means. The 2.5th and
97.5th percentiles of the replicate distribution give each model's confidence
interval. The board publishes intervals for efficacy and the CLEAR composite.

\paragraph{Paired normalization.} Because the CLEAR composite is field-relative
(min-max normalized across the evaluated models), its uncertainty must include
normalization variance, not only the variance of each raw dimension. The
\emph{same} resampled set of scenario IDs is applied to every model within a
replicate (a paired bootstrap), and the full normalization plus weighted composite
is recomputed inside each replicate across all models. This keeps the field
consistent per replicate and propagates normalization variance into the CLEAR
interval. Pairing on the scenario set is also what makes model-versus-model
difference intervals valid, which is what ``is A really better than B'' requires.

\paragraph{Macro-averaged scores.} The headline efficacy is the micro average
over scenarios, which lets a frequent-easy scenario category drown a rare-hard
one. The board therefore also publishes macro-averaged efficacy — the unweighted
mean of per-category (and per-domain) means — alongside the micro value. Its
interval comes from a category-stratified variant of the same paired bootstrap
(scenarios resampled with replacement \emph{within} each category, then category
means averaged), with the same $B$, seed, and run-collapsing conventions.

\subsection{Rank bands}

Orderings are only meaningful where intervals separate. Models are clustered into
rank bands by CLEAR-CI overlap with a greedy walk over models sorted by CLEAR
descending: the current band's leader is the highest unbanded model, every
following model whose CLEAR interval overlaps the leader's joins the band, and the
first non-overlapping model starts a new band. Models within a band are not
statistically distinguishable and are marked with a tied rank. This is the
Arena-Hard ``separability with confidence'' idea applied to absolute
scores~\cite{arenahard2024}: the honest way to present a roster where many gaps are
noise. The leaderboard JSON carries a top-level note stating the scenario count and
the band caveat.

\subsection{Publish minimum}

The board does not publish at all when any evaluated domain has fewer than 30
scenarios. Below that, the bootstrap intervals on every dimension overlap almost
completely, so any ranking is indistinguishable from chance. A publish gate
enforces this alongside a run-completeness check; a partial flag downgrades it to a
warning for deliberate previews. Edge cases are handled without crashing: a single
scenario yields a degenerate zero-width interval, and a single-model run skips
cross-model normalization so its CLEAR interval mirrors its efficacy interval.

\subsection{Inter-judge reliability: Krippendorff's $\alpha$}

A raw within-0.2 agreement rate is not chance-corrected, not comparable across
score distributions, and uses an arbitrary threshold. Krippendorff's $\alpha$ is the
field-standard chance-corrected metric for three or more raters on ordinal or
interval labels (Cohen's $\kappa$ is for two raters and assumes nominal categories).
$\alpha$ is computed at the interval level (squared-difference distance, appropriate
for continuous 0--1 scores) over the published per-judge columns: each result row
is a unit, each judge a rater, and a parse-failed judge is a missing value, which
$\alpha$ handles natively (it uses only units with two or more present scores).
$\alpha$ is published per dimension both overall and per model. It is implemented
in-repo as a small test-validated function checked against worked examples from
Krippendorff's reference material~\cite{krippendorff2011}, rather than adding a
dependency, and is null when undefined. The within-0.2 rate is kept as a secondary
human-readable readout.

\subsection{Length-bias regression}

LLM judges can favor verbose output, the AlpacaEval length-bias
lesson~\cite{lengthcontrolled2024}. Rather than assert immunity, the benchmark
measures it. At aggregation, an ordinary-least-squares regression fits each judge
dimension's consensus on the agent's \texttt{output\_tokens} across all rows. The
bias is a property of the judge panel, not of any one model, so the regression
pools all rows; the deterministic state score is judge-independent and excluded.
For each dimension the board publishes the slope, intercept, $R^2$, the slope's
standard error and $t$-statistic, the sample size, and a significance flag
($|t| > 1.96$). A materially positive, significant slope is direct evidence that
longer outputs earn higher judge scores independent of correctness, and is the
trigger for adding an explicit length-neutrality line to the rubrics. This is the
single highest-leverage, lowest-effort robustness check available, because the
output-length data already exists for the cost dimension.
